\chapter{训练顺序的计算动力学：时间索引、路径依赖与持续学习}
\label{chp:training_paradigm_original_sin}

训练过程不仅决定模型参数落在何处，也决定哪些证据能够相遇、哪些更新能够保留，以及系统是否具备区分观察时间、事件时间和自身更新历史的接口。但这不意味着随机洗牌必然抹除因果，更不意味着给表示增加一个时间坐标就会自动产生连续自我。我们把原章关于“维度坍缩”和“热力学粉碎”的断言，重建为可比较的训练协议：明确样本来源、事件顺序、参数版本、外部记忆与评价任务，测量不同排序和更新机制对泛化、校准、遗忘、因果干预与身份连续性的影响。本章据此讨论时间信息、非对易更新、因果证据和持续学习，并始终区分数据中的历史、模型的计算状态与部署主体的经历。

\section{时间不是额外坐标：事件索引、版本状态与可查询历史}
\label{sec:dimensional_deficit_time}

我们先区分四类经常被“时间维度”一词混合的对象。第一类是事件时间$\tau_e$，表示被描述或观测的事情何时发生，单位可以是秒、日或历史年代；第二类是记录时间$\tau_r$，表示数据何时进入日志、档案或训练集合；第三类是训练步$k\in\mathbb{N}$，它是优化器执行更新的离散序号，不等同于现实时间；第四类是会话或主体历史$H_k$，记录系统在版本$k$之前接收过的观测、采取过的行动、获得的反馈和被授权保留的承诺。文本中的“昨天”若缺少参照时刻，确实不能唯一定位事件；但带有日期的语料也不会仅凭日期字段变成主体亲历。事件索引属于数据语义，训练步属于算法执行，历史记录属于运行状态。三者可通过来源身份、版本号与查询接口绑定，却不能压成一个所谓$+1$维坐标。

设第$k$步训练状态为
\[
Z_k=(\theta_k,m_k,M_k,C_k),
\]
其中$\theta_k\in\mathbb{R}^d$是模型参数，$m_k$是优化器动量等辅助量，$M_k$是可写外部记忆，$C_k$是数据来源、访问权限和版本承诺组成的控制状态。一个带顺序的样本批$B_k=\{(x_i,y_i,\tau_{e,i},\tau_{r,i},s_i)\}$还包含来源$s_i$。离散更新写成
\[
Z_{k+1}=\mathcal{U}_{\eta_k}(Z_k,B_k,q_k,\xi_k),
\]
其中$\eta_k$是步长，$q_k$是当步目标和约束，$\xi_k$表示采样与数值噪声。若数据分布、目标、权限或模型结构随$k$变化，就必须显式进入$\mathcal{U}$；只写一个固定自治方程会把课程变化、反馈学习和部署修订误当成相同过程。这个定义说明，训练天然具有算法时间和版本历史，即使输入没有现实时间戳，参数也不会处于“永恒当下”。真正的问题是：更新后哪些历史差异仍能由任务相关查询读出，哪些被参数共享和压缩不可逆地合并。

时间信息的必要程度取决于任务。语言建模可以从句序、时态和叙事结构学习大量相对顺序；知识问答若涉及会变化的事实，则需要来源日期、有效期和检索时刻；机器人控制还要记录传感器时延、动作发出时间与环境响应；长期助理则必须区分“用户去年表达过的偏好”“本次会话的临时要求”和“系统当前获准执行的规则”。因此，我们不把缺少绝对空间坐标视为所有语言模型的几何残缺。相反，我们为每类任务规定最小时间契约$T_q$：允许的时钟、精度、参照系、缺失值语义以及冲突解决方法。一个历史文本可以只有世纪级精度，一个控制回路可能要求毫秒级同步；把两者统一到同一绝对坐标既无必要，也可能制造虚假精确性。

“知道”与“经历过”也必须通过数据来源和因果参与来区分。令$e$为事件，$a$为系统动作，$o$为后续观测。系统知道事件，可以来自记录$r(e)$的检索或参数中的统计编码；系统经历事件，则至少要求存在同一持续身份下的闭环链$H_k\supset(o_{k-1},a_k,o_k)$，并能回答该动作如何改变了后续可观测状态。即便如此，我们也只能主张它具备来源记忆、行动归因和历史连续性等功能，不由这些工程指标推出主观体验。反过来，没有在线行动的模型仍可能借助带来源的外部记忆准确追踪事实版本。把“亲历”设为所有知识的前提，会错误排除文献证据、社会学习和跨代传承。

可查询历史需要结构网络与分布表示协作。事件日志形成局部结构网络：节点保存事件、主体、对象和版本身份，边表示先后、因果候选、授权与引用；模型参数和活动状态则提供跨样本的分布式相似性与预测。检索时，任务身份把查询绑定到相关日志片段，读出接口把证据送入当前状态，控制器再决定是否更新参数、外部记忆或仅生成一次性回答。AgentOS可以把这一分工实现为$h(t)$、外部记忆$M_e$、Boundary与Offloading，但这些只是工程实例；其他系统可以采用数据库、事件溯源、可微记忆或群体档案，只要满足身份、权限、时序与可回滚性要求。

经验上，我们用对照协议检验时间表示，而不是从伪黎曼度量推出结论。对同一数据集合构造四组训练：完全随机排序、保留局部文档顺序、按事件时间课程化、按因果先决关系排序；控制样本数量、优化器、计算预算和评价时刻。评价至少包括过时事实更新、时间关系问答、反事实干预、长程承诺保持、来源追溯和灾难性遗忘。若课程化只改善早期收敛而最终性能相同，它是优化效应；若来源时间只改善会变事实而不改善一般推理，它是检索校准效应；若带闭环历史的系统在行动归因和承诺一致性上稳定胜出，才支持持续身份机制的主张。由此，原先的“缺失$+1$维”被替换成更严格的问题：给定任务要求，系统是否保存了足够的时间索引、版本状态与可查询历史，以及这些结构能否在分布变化和权限约束下继续工作。
